\section{Related work in more detail}
\label{app:related}

\paragraph{Auditory attention decoding.}
Brain activity follows the envelope of the speech a listener attends to
\citep{ding2012,lalor2009}. Standard decoders either rebuild the envelope from
EEG and correlate it with each candidate \citep{osullivan2015}, or learn the
match directly \citep{accou2021,monesi2020,accou2023,puffay2023}. A separate line
decodes the \emph{direction} of attention without any audio
\citep{vandecappelle2021}. The work closest to ours is \citet{rotaru2024}, who
showed that direction decoding on several public datasets is driven largely by
eye movements and by regularities in the trial design, not by attention. We bring
the same concern to the matching task and to several recorded streams at once,
and we make the check part of how models are chosen and reported rather than a
test run afterwards.

\paragraph{Shortcuts in datasets.}
Models use whatever solves the task \citep{geirhos2020}, and benchmarks often
contain cues that make the intended input unnecessary
\citep{torralba2011,gururangan2018,goyal2017vqa}. The usual fixes are to rebalance
the data \citep{goyal2017vqa} or to train against a branch that sees only the
shortcut \citep{cadene2019rubi}. \merit{} does both --- it checks the candidate
sets and it trains an audio-only adversary \citep{ganin2015} --- and adds one thing
specific to the matching task: a scoring function in which the shortcut is
impossible rather than discouraged (Section~\ref{sec:coupling}).

\paragraph{Balancing modalities.}
When modalities are trained together, the easiest one dominates, and the fused
network can end up worse than its best single part \citep{wang2020gradblend}.
This has been called modality competition \citep{huang2022competition} and a
greedy learning process \citep{wu2022greedy}, and it is usually treated by
evening out \emph{training} --- adjusting gradients as training runs
\citep{peng2022ogm}, blending losses \citep{wang2020gradblend}, or dropping
modalities at random. \merit{} looks at something different. Those methods balance
how fast each modality learns; a scientific claim instead rests on how much each
modality affects the decision \emph{at test time}, which is what
Eq.~\ref{eq:shapley} reports. The two are not the same: in Table~\ref{tab:hinge},
turning modality dropout on or off changes the EEG's credit very little, while the
EEG's credit is still far below what it earns alone.

\paragraph{Attribution.}
Permutation importance \citep{breiman2001,fisher2019} measures how much a model
depends on an input by breaking that input's link to the rest of the example.
Shapley values \citep{shapley1953,lundberg2017} give the only way to split credit
that adds up exactly, and SAGE \citep{covert2020} uses them for importance measured
by performance over a whole dataset. We apply the same idea to whole streams,
where there are few enough subsets ($2^M$) to compute it exactly;
\citet{hessel2020} make a related point for multimodal models. What is new is not
the measurement but how it is used: to choose checkpoints, and to report the
brain's share alongside every accuracy. We also tried adding it to the training
loss, and report that this does not help (Section~\ref{sec:res-hinge}).

\paragraph{Components.}
The contrastive term is InfoNCE \citep{oord2018} in the two-way arrangement of
CLIP \citep{radford2021clip}, with negatives from the same listener; the
anti-collapse term follows VICReg \citep{bardes2022vicreg}; and the foveal stream
is a frozen V-JEPA\,2 encoder \citep{assran2025vjepa2}.
